\documentclass[10pt, titlepage]{article}

\usepackage[margin=1in]{geometry}
\usepackage{fancyhdr}
\usepackage{amsfonts}
\usepackage{amsmath}
\usepackage{amssymb}
\usepackage{booktabs}
\usepackage{tikz}
\usepackage{float}
% \usepackage{tikz}
% \usepackage{pgfplots}


\makeatletter
\setlength{\@fptop}{0pt}
\makeatother

\title{\normalsize PLUS Session Worksheet \\[1em] \Large \textit{MATH 1554:} Linear Algebra}
\author{\textit{Exam 3 Review}}
\date{Summer 2023}

% ------- Set up header -------
\renewcommand{\sectionmark}[1]{\markright{#1}} % from https://tex.stackexchange.com/a/317393/227665
\pagestyle{fancy}
\fancyhead[L]{Summer 2023}
\fancyhead[C]{\textit{\nouppercase{\rightmark}}}
% \fancyhead[C]{\textit{(6.1), (6.2), \& (6.3)}}
\fancyhead[R]{\textbf{\thepage}}
% -------------------------------------------------------

% ------- Set up question environment -------
\newcounter{question}
\setcounter{question}{1}
\newenvironment{newquestion}
{
	\noindent\textbf{Question \thequestion}
}
{
\stepcounter{question}
\par
}
% -------------------------------------------------------


% ------- Set up math stuff -------
\renewcommand{\vec}[1]{\mathbf{#1}}
\DeclareMathOperator{\vspan}{Span}
\DeclareMathOperator{\row}{row}
\DeclareMathOperator{\col}{col}
\DeclareMathOperator{\colspace}{Col}
\DeclareMathOperator{\rowspace}{Row}
\DeclareMathOperator{\nullspace}{Nul}
\DeclareMathOperator{\rank}{rank}
\DeclareMathOperator{\proj}{proj}
% -------------------------------------------------------

% ------- Set up circle -------
\newcommand{\bubble}{\tikz[baseline=-0.625ex]\draw[radius=0.5em, fill=white] (0,0) circle;}
\newcommand{\dbubble}{\tikz[baseline=-0.625ex]\draw[radius=0.5em, fill=black] (0,0) circle;}
% -------------------------------------------------------

\begin{document}
\maketitle

\section{\centering T \& F}
	\begin{newquestion}
		For the following statements, answer \textit{true} if the statement is always true, and \textit{false} otherwise.
		\begin{table}[!h]
			\centering
			\begin{tabular}{c@{\hspace{2.5em}}cp{0.75\textwidth}}
				true & false & \\
				\specialrule{\lightrulewidth}{0.25em}{1.25em}
				\bubble & \bubble & For any scalar \(c \in \mathbb{R}\), \(\|c\vec{v}\| = c\|\vec{v}\|\). \\[0.75em]
				\bubble & \bubble & If \(\vec{u} \cdot \vec{v} = 0\), then \(\|\vec{v} - \vec{u}\| = \|\vec{v} + \vec{u}\|.\)\\[0.75em]
				\bubble & \bubble & If \(\vec{y} \in \mathbb{R}^n\) is orthogonal to some \(\vec{z} \in W\), 
							      where \(W\) is a subspace of \(\mathbb{R}^n\), then \(\vec{y} \in W^\perp\). \\[0.75em]
				\bubble & \bubble & If \(A \in \mathbb{R}^{10 \times 12}\) and \(\rank A = 6\), then \(\dim (\nullspace A^T) = 6\).\\[0.75em]
				\bubble & \bubble & For any \(A \in \mathbb{R}^{m \times n}\), \(\nullspace A^TA = \nullspace A\). \\[0.75em]
				\bubble & \bubble & If \(A\) and \(B\) are orthogonal matrices of the same size, then so is \(AB\). \\[0.75em]
				\bubble & \bubble & The Gram-Schmidt process on a set of vectors in \(\mathbb{R}^n\) always yields a basis for the subspace formed
							      by their span. \\[0.75em]
				\bubble & \bubble & An invertible matrix \(A \in \mathbb{R}^n\) is diagonalizable if and only if it's inverse is. \\[0.75em]
				\bubble & \bubble & If \(\vec{y} \in \colspace B\) and \(B\) has orthonormal columns, then \(BB^T\vec{y} = \vec{y}\). \\[0.75em]
				\bottomrule
			\end{tabular}
		\end{table}
	\end{newquestion}

\section{\centering Possible/Impossible}

	\begin{newquestion}
		Answer the following statements with \textit{possible} if the statement \emph{could} be true.
		\begin{table}[!h]
			\centering
			\begin{tabular}{c@{\hspace{2.5em}}cp{0.75\textwidth}}
				possible & impossible & \\
				\specialrule{\lightrulewidth}{0.25em}{1.25em}
				\bubble & \bubble & An orthogonal matrix \(U\) with \(\lambda = 0\) as an eigenvalue. \\[0.75em]
				\bubble & \bubble & Two vectors \(\vec{v}_1 \ne \vec{v}_2\) in \(\mathbb{R}^3\) with 
							      \(\vec{v}_1, \vec{v}_2 \in W^\perp\) and \(\vec{y} \in W\) where 
							      \(\|\vec{y} - \vec{v}_1\| = \|\vec{y} - \vec{v}_2\|\) and \(W\) is a 2-dimensional subspace of 
							       \(\mathbb{R}^3\).\\[0.75em]
				\bubble & \bubble & \(\vec{v} \ne \vec{0}\) and \(\vec{u} \ne \vec{0}\) but 
							      \(\proj_{\vec{v}} \vec{u} = \vec{0} = \proj_{\vec{u}} \vec{v}\). \\[0.75em]
				\bubble & \bubble & Two vectors that form a linearly dependent set but also an orthogonal set. \\[0.75em]
				\bubble & \bubble & A \(QR\) factorization of \(A\) where \(\colspace A \ne \colspace Q\). \\[0.75em]
				\bubble & \bubble & A \(4 \times 4\) diagonalizable matrix with 2 distinct eigenvalues of geometric multiplicities 2, 1 and a total algebraic
							      multiplicity of 4.\\[0.75em]
				\bottomrule
			\end{tabular}
		\end{table}
	\end{newquestion}

\section{\centering Short Answer}
	\begin{enumerate}
		\item Let \(A = \small \begin{pmatrix}1/\sqrt{2} & 0 & -1/\sqrt{2}\\ 0 & 1 & 0\\ 1/\sqrt{2} & 0 & 1/\sqrt{2}\end{pmatrix}\) and 
		\(\vec{v} = \begin{pmatrix}1 \\ 0 \\ 3\end{pmatrix}\). Find \(\|A\vec{v}\|\).
		\item If \(W = \{\vec{x} \in \mathbb{R}^4 \mid x_1 + x_2 = x_3 \ \text{and} \ x_4 = 0\}\), then what is \(\dim W^\perp\)? \(\dim W\)?
		\item Find a 2-by-3 matrix \(A\) such that \((\colspace A)^\perp = \vspan \biggl\{ \begin{bmatrix}2 \\ 2\end{bmatrix}\biggr\}\).
		\item If \(\|\vec{v}\| = \|\vec{u}\| = 1\) and \(\vec{v} \cdot \vec{u} = 0\), then what is \(\|2\vec{v} + 3\vec{u}\|\)?
		\item If \(A \in \mathbb{R}^{16 \times 15}\) and \(\dim (\nullspace A^T) = 12\), what is \(\dim((\nullspace A)^\perp) + \dim (\colspace A)\)?
		\item If \(\vec{u}_1 = \begin{bmatrix}1 \\ 0 \\ 1\end{bmatrix}, \vec{u}_2 = \begin{bmatrix}1 \\ 1 \\ -1\end{bmatrix}, 
			      \vec{u}_3 = \begin{bmatrix}-1 \\ 2 \\ 1\end{bmatrix}\), and \(W = \vspan \{\vec{u}_1, \vec{u}_2\}\):
			\begin{enumerate}
				\item What is the closest vector in \(W\) to \(\begin{bmatrix}1 \\ 5 \\ 1\end{bmatrix}\)?
				\item Find the coordinate vector of \(\begin{bmatrix}3 \\ 1 \\ 1\end{bmatrix}\) under the basis formed by 
				\(\mathcal{B} = \{\vec{u}_1, \vec{u}_2\}\).
				\item Find a vector \(\vec{x} \in \mathbb{R}^3\) such that:
					\begin{enumerate}
						\item \(\vec{x}\) is orthogonal to \(\vec{u}_3\).
						\item \(\proj_{\mathcal{L}_1} \vec{x} = 5\vec{u}_1,\) where \(\mathcal{L}_1 = \vspan \{\vec{u}_1\}\)
						\item \(\proj_{\mathcal{L}_2} \vec{x} = 3\vec{u}_2,\) where \(\mathcal{L}_2 = \vspan\{\vec{u}_2\}\)
					\end{enumerate} 
			\end{enumerate}
		\item Find the value(s) for the parameters \(\alpha, \beta\) that'll best fit the curve \[y = \alpha x^2 + \beta\] to the data points 
			\((-1, 1), (1, 3), (1, 2), (1, 1)\).
		\item Use the \(QR\) decomposition of the matrix \(A = QR = \frac{1}{\sqrt{2}}\begin{bmatrix}1 & 1 \\ -1 & 1\end{bmatrix}\begin{bmatrix}2 & 2 \\0 & 1\end{bmatrix}\) to find the least square solution of \(A\vec{x} = \begin{bmatrix}2\sqrt{2} \\ \sqrt{2}\end{bmatrix}\).

		\item Find the least-square solution for the system
			\(\begin{bmatrix}1 & 1 \\ 0 & 0 \\ -1 & 1\\ 0 & 0\end{bmatrix}\vec{x} =  \begin{bmatrix}1 \\ 2 \\ 4 \\ 1\end{bmatrix}\).
	\end{enumerate}

\newpage
\fancyhead[C]{\textit{Solutions}}

\section{\centering Solutions to T \& F Questions}
	\begin{table}[!h]
		\centering
		\begin{tabular}{c@{\hspace{2.5em}}cp{0.75\textwidth}}
			true & false & \\
			\specialrule{\lightrulewidth}{0.25em}{1.25em}
			\bubble & \dbubble & For any scalar \(c \in \mathbb{R}\), \(\|c\vec{v}\| = c\|\vec{v}\|\). \\[0.75em]
			\dbubble & \bubble & If \(\vec{u} \cdot \vec{v} = 0\), then \(\|\vec{v} - \vec{u}\| = \|\vec{v} + \vec{u}\|.\)\\[0.75em]
			\bubble & \dbubble & If \(\vec{y} \in \mathbb{R}^n\) is orthogonal to some \(\vec{z} \in W\), 
						       where \(W\) is a subspace of \(\mathbb{R}^n\), then \(\vec{y} \in W^\perp\). \\[0.75em]
			\bubble & \dbubble & If \(A \in \mathbb{R}^{10 \times 12}\) and \(\rank A = 6\), then \(\dim (\nullspace A^T) = 6\).\\[0.75em]
			\dbubble & \bubble & For any \(A \in \mathbb{R}^{m \times n}\), \(\nullspace A^TA = \nullspace A\). \\[0.75em]
			\dbubble & \bubble & If \(A\) and \(B\) are orthogonal matrices of the same size, then so is \(AB\). \\[0.75em]
			\bubble & \dbubble & The Gram-Schmidt process on a set of vectors in \(\mathbb{R}^n\) always yields a basis for the subspace formed
						      by their span. \\[0.75em]
			\dbubble & \bubble & An invertible matrix \(A \in \mathbb{R}^n\) is diagonalizable if and only if it's inverse is. \\[0.75em]
			\dbubble & \bubble & If \(\vec{y} \in \colspace B\) and \(B\) has orthonormal columns, then \(BB^T\vec{y} = \vec{y}\). \\[0.75em]
			\bottomrule
		\end{tabular}
	\end{table}

\subsection{Justification to Solutions To T \& F Questions}
	\paragraph{Justification to \textit{Q1}} It should be \(\|c\vec{v}\| = \lvert c \rvert \|\vec{v}\|\), since the norm has to always be positive.
	
	\paragraph{Justification to \textit{Q2}} You can draw a picture and convince yourself. Here is how you'd show it algebraically.
	\begin{align*}
		\|\vec{v} - \vec{u}\|^2 &= (\vec{v} - \vec{u})^T(\vec{v} - \vec{u}) \\
					           &= \vec{v}^T\vec{v} - \vec{v}^T\vec{u} - \vec{u}^T\vec{v} + \vec{u}^T\vec{u} \\
	                                           &= \|\vec{v}\|^2 + \|\vec{u}\|^2
	\end{align*}
	\begin{align*}
		\|\vec{v} + \vec{u}\|^2 &= (\vec{v} + \vec{u})^T(\vec{v} + \vec{u}) \\
				           &= \vec{v}^T\vec{v} + \vec{v}^T\vec{u} + \vec{u}^T\vec{v} + \vec{u}^T\vec{u} \\
                                           &= \|\vec{v}\|^2 + \|\vec{u}\|^2
	\end{align*}
	Since the norm always yields a positive number, we know that 
	\(\|\vec{v} - \vec{u}\|^2 = \|\vec{v} + \vec{u}\|^2 \implies \|\vec{v} - \vec{u}\| = \|\vec{v} + \vec{u}\|\). Therefore, 
	\(\|\vec{v} - \vec{u}\| = \|\vec{v} + \vec{u}\|\) if \(\vec{v} \perp \vec{u}\).
	
	\paragraph{Justification to \textit{Q3}} By definition, \(\vec{y}\) has to be orthogonal to \emph{every} vector in \(W\). Just picture two vectors that are in a
			plane in \(\mathbb{R}^3\). They \emph{might} be orthogonal to each other, but they're both in \(W\).
	
	\paragraph{Justification to \textit{Q4}} For any subspace \(W\) of \(\mathbb{R}^n\), \(\dim W + \dim W^\perp = n\). In this case, we know that 
			\((\colspace A)^\perp = \nullspace A^T\). Since \(\rank A = \dim(\colspace A) = 6\), and \(\dim(\colspace A) + \dim(\nullspace A^T) = 10,\)
			\(\dim(\nullspace A^T) = 4 \ne 6\).\footnote{The reason they add up to 10 rather than 12 is because they are both subspaces of 
			\(\mathbb{R}^{10}\).}
	
	\paragraph{Justification to \textit{Q5}} To show that two sets are equal, you have to show that they are subsets of each other. Basically, in this case,
			you have to show that every vector in \(\nullspace A\) also exists in \(\nullspace A^TA\), and that every vector in \(\nullspace A^TA\) 
			\emph{also} is in \(\nullspace A\).
			Let's show that \(\nullspace A \subseteq \nullspace A^TA\) first. Choose an arbitrary vector \(\vec{x} \in \nullspace A.\) The vector satifies the
			equation
			\[A\vec{x} = \vec{0} \implies A^TA\vec{x} = A^T \vec{0} \implies A^TA\vec{x} = \vec{0}.\] This arbitrary \(\vec{x}\) is therefore also in
			\(\nullspace A^TA\). Therefore, every vector in \(\nullspace A\) is also in \(\nullspace A^TA\). 

			Now, we go ther other way. Pick an arbitrary vector \(\vec{z} \in \nullspace A^TA.\) This vector satisfies the equation
			\[A^TA \vec{z} = \vec{0} \implies \vec{z}^T A^TA \vec{z} = \vec{0} \implies (A\vec{z}) \cdot (A\vec{z}) = \vec{0}.\]
			It is important to realize that \(\vec{v} \cdot \vec{v} = 0 \iff \vec{v} = \vec{0}.\) So, \(A\vec{z} = \vec{0},\) meaning that 
			\(\vec{z} \in \nullspace A\) as well. Note that \(\vec{z}\) could be non-zero as well if the nullspace of \(A^TA\) is not trivial. Therefore, we
			showed that \(\nullspace A^TA \subseteq \nullspace A\).

			In conclusion, \(\nullspace A^TA = \nullspace A\)!

	\paragraph{Justification to \textit{Q6}} \((AB)(AB)^T = A(BB^T)A^T = AA^T = I.\) It is trivial to show that \(AB\) is left invertible as well with the same
			inverse. So \(AB\) is orthogonal as well.

	\paragraph{Justification to \textit{Q7}} Apply Gram-Schmidt on three vectors in \(\mathbb{R}^2.\) Recall that a basis has to be linearly independent as one
			of the conditions.

	\paragraph{Justification to \textit{Q8}} If \(A\) is diagonalizible and invertible, then \[A = PDP^{-1} \implies A^{-1} = (PDP^{-1})^{-1} = PD^{-1}P^{-1}.\]
			therefore, the inverse of a diagonalizble invertible matrix is diagonalizable.

	\paragraph{Justification to \textit{Q9}} If \(B\)'s columns form an orthonormal set, then \(BB^T\) is the projection matrix. As in\footnote{You can multiply the matrices out and you'd get the projection formula we're used to.}
			\[\proj_{\colspace B}(\vec{x}) = BB^T\vec{x}.\] If \(\vec{y} \in \colspace B\), then its projection onto the columnspace should intuitvely be itself.

\section{Solutions to Possible/Impossible Questions}
	\begin{table}[H]
		\centering
		\begin{tabular}{c@{\hspace{2.5em}}cp{0.75\textwidth}}
			possible & impossible & \\
			\specialrule{\lightrulewidth}{0.25em}{1.25em}
			\bubble & \dbubble & An orthogonal matrix \(U\) with \(\lambda = 0\) as an eigenvalue. \\[0.75em]
			\dbubble & \bubble & Two vectors \(\vec{v}_1 \ne \vec{v}_2\) in \(\mathbb{R}^3\) with 
						      \(\vec{v}_1, \vec{v}_2 \in W^\perp\) and \(\vec{y} \in W\) where 
						      \(\|\vec{y} - \vec{v}_1\| = \|\vec{y} - \vec{v}_2\|\) and \(W\) is a 2-dimensional subspace of 
						       \(\mathbb{R}^3\).\\[0.75em]
			\dbubble & \bubble & \(\vec{v} \ne \vec{0}\) and \(\vec{u} \ne \vec{0}\) but 
						      \(\proj_{\vec{v}} \vec{u} = \vec{0} = \proj_{\vec{u}} \vec{v}\). \\[0.75em]
			\dbubble & \bubble & Two vectors that form a linearly dependent set but also an orthogonal set. \\[0.75em]
			\bubble & \dbubble & A \(QR\) factorization of \(A\) where \(\colspace A \ne \colspace Q\). \\[0.75em]
			\bubble & \dbubble & A \(4 \times 4\) diagonalizable matrix with 2 distinct eigenvalues of geometric multiplicities 2, 1 and a total algebraic
							multiplicity of 4.\\[0.75em]
			\bottomrule
		\end{tabular}
	\end{table}

\subsection{Justification to Possible/Impossible Questions}
	\paragraph{Justification to \textit{Q1}} Every orthogonal matrix is invertible, with it's inverse being its transpose.
	
	\paragraph{Justification to \textit{Q2}} This is just an extension of the argument from (\textbf{4.1}, \textit{Justification to Q2}) but extended to a higher
			dimension. Draw a picture and everything should be clear.

	\paragraph{Justification to \textit{Q3}} Since orthogonal vectors don't ``share anything in common" (vectors in an orthogonal set do not ``overlap"), 
			they're projection into the subspace formed by each other's span should yield the zero vector. This should be consistent with the projection 
			formula we know.

	\paragraph{Justification to \textit{Q4}} Note that the zero vector could be in the set.

	\paragraph{Justification to \textit{Q5}} This is one of the promises of the Gram-Schmidt process. Also, note that in \(QR\) factorization, the matrix
			\(R\) is upper triangular, and it's job is to take linear combinations of the columns of \(Q\) to yield \(A\). This means that there column space is
			also the same.

	\paragraph{Justification to \textit{Q6}} This is impossible. The geometric multiplicity has to add up to 4 in order to be diagonalizable in this case.
\subsection{Solutions to Short Answer}
	\begin{enumerate}
		\item Since \(A\) is an orthogonal matrix, then \(\|A\vec{v}\| = \|\vec{v}\| = \sqrt{\vec{v} \cdot \vec{v}} = \sqrt{10}\).
		\item Note that \(W = \{\vec{x} \in \mathbb{R}^4 \mid x_1 + x_2 = x_3 \ \text{and} \ x_4 = 0\} = \{\vec{x} \in \mathbb{R}^4 \mid x_1 + x_2 - x_3 = 0 
			\ \text{and} \ x_4 = 0\}\). This indicates that \(W\) is the solution set to the homogenous system 
			\[\begin{aligned}
				x_1 + x_2 - x_3 \hphantom{+ x_4}  &= 0 \\
				                                            x_4   &= 0
			\end{aligned}\]
			We immediately see that the system has 2 free variables, so \(\dim W = 2\). Since we're working in \(\mathbb{R}^4, \) 
			\(\dim W + \dim W^\perp = 4 \implies \dim W^\perp = 2.\)
		\item Multiple answers. One such answer is \[\begin{bmatrix}1 & 1 & 1 \\ -1 & -1 & -1\end{bmatrix}.\] The most general answer is in the form 
			\[\begin{bmatrix}a & ca & da \\ -a & -ca & -da\end{bmatrix},\] where \(c, d \in \mathbb{R}.\) Note that the question is simply asking for
			a matrix whose columns are orthogonal to \(\begin{bmatrix}2 \\ 2\end{bmatrix}\). We know that the columnspace is of dimension 1 because
			of the fact that \(\dim (\colspace A)^\perp = 1\). This means, you only need to find one vector orthogonal to 
			\(\begin{bmatrix}2 \\ 2\end{bmatrix}\) and the rest has to be a scalar multiple of that.
		\item \(\|2\vec{v} + 3\vec{u}\|^2 = 4\|\vec{v}\|^2 + 9\|\vec{u}\|^2 = 13.\) We've used the Pythagorean Theorem. 
			So \(\|2\vec{v} + 3\vec{u}\| = \sqrt{13}\).
		\item \((\nullspace A)^\perp = \rowspace A,\) and \(\dim(\rowspace A) = \dim(\colspace A) = 4.\) So the sum is \textbf{8}. 
		\item Answers are below.
			\begin{enumerate}
				\item The closest vector to in \(W\) to \(\vec{v} = \begin{bmatrix}1 \\ 5 \\ 1\end{bmatrix}\) is \(\proj_W \vec{v}.\)
				In this case, that's 
				\[\proj_W \vec{v}
				  = \frac{\vec{v} \cdot \vec{u}_1}{\vec{u}_1 \cdot \vec{u}_1}\vec{u}_1 
				   + \frac{\vec{v} \cdot \vec{u}_2}{\vec{u}_2 \cdot \vec{u}_2}\vec{u}_2 = \begin{bmatrix}8/3 \\ 5/3 \\ -2/3\end{bmatrix}\]
				\item Since the vector is already in \(W\), finding it's projection amounts to finding a linear combination of \(\{\vec{u}_1, \vec{u}_2\}\)
					that'll yield \(\vec{v} = \begin{bmatrix}3 \\ 1 \\ 1\end{bmatrix}\). That is 
					\[\proj_W \vec{v} = \frac{\vec{u}_1 \cdot \vec{v}}{\vec{u}_1 \cdot \vec{u}_1}\vec{u}_1 + \frac{\vec{u}_2 \cdot \vec{v}}{\vec{u}_2 \cdot \vec{u}_2}\vec{u}_2 = \begin{bmatrix}\vec{u}_1 & \vec{u}_2\end{bmatrix}\begin{bmatrix}\frac{\vec{u}_1 \cdot \vec{v}}{\vec{u}_1 \cdot \vec{u}_1}\\[0.5em] \frac{\vec{u}_2 \cdot \vec{v}}{\vec{u}_2 \cdot \vec{u}_2}\end{bmatrix}.\] The final vector is the coordinate of \(\vec{v}\) under the new basis. 
					Computing the values will show that the coordinate is \(\begin{bmatrix}2 \\ 1\end{bmatrix}\).
				\item Note that \(\{\vec{u}_1, \vec{u}_2, \vec{u}_3\}\) forms an orthogonal set. Therefore, if \(W =  \vspan\{\vec{u}_1, \vec{u}_2\}\),
					 then \(W^\perp = \vspan\{\vec{u}_3\}\).\footnote{This is shown in the textbook, and it's easy to verify for yourself.} Since we want
					 the vector to be orthogonal to \(\vec{u}_3\), it should live in \(W\). Since the vector already lives in \(W\), then the weights of the 
					projection are just its new coordinates under the basis formed by \(\mathcal{B} = \{\vec{u}_1, \vec{u}_2\}.\) (If you tell me the
					projection, you're also telling me the coordinates in this case!) This alone is enough to tell us that \(\vec{x} = 5\vec{u}_1 + 3\vec{u}_2.\)
				\item Follow the typical procedure for solving these types of problems.
					\[\begin{aligned}
						1 &= \alpha (-1)^2 + \beta &&\implies 1 = \alpha + \beta \\
						3 &= \alpha (1)^2  + \beta &&\implies 3 = \alpha + \beta \\
						2 &= \alpha (1)^2  + \beta &&\implies 2 = \alpha + \beta \\
						1 &= \alpha (1)^2  + \beta &&\implies 1 = \alpha + \beta \\
					\end{aligned}\]
					Something weird is happening here. The matrix equation associated with this system looks like 
					\[\begin{bmatrix}1 & 1 \\ 1 & 1 \\ 1 & 1 \\ 1 & 1\end{bmatrix}\begin{bmatrix}\alpha \\ \beta\end{bmatrix} 
					  = \begin{bmatrix}1 \\ 3 \\ 2 \\ 1\end{bmatrix}\]
					\textbf{The columns do not form a linearly independent set, so we'll have an infinite number of least-square solutions for 
					\(\alpha, \beta\) where they all yield the same amount of error}. Moving on,
					\[A^TA = \begin{bmatrix}1 & 1 & 1 & 1 \\ 1 & 1 & 1 & 1\end{bmatrix}
					\begin{bmatrix}1 & 1 \\ 1 & 1 \\ 1 & 1 \\ 1 & 1\end{bmatrix} = \begin{bmatrix}4 & 4 \\ 4 & 4\end{bmatrix};\]

					\[A^T\vec{b} = \begin{bmatrix}1 & 1 & 1 & 1 \\ 1 & 1 & 1 & 1\end{bmatrix}\begin{bmatrix}1 \\ 3 \\ 2 \\ 1\end{bmatrix}
					 = \begin{bmatrix}7 \\ 7\end{bmatrix}.\]

					Now, we solve the normal equations \(A^TA\vec{\hat{x}} = A^T\vec{b}.\) 
					\[\left(\!\!\begin{array}{cc|c}4 & 4 & 7 \\ 4 & 4 & 7\end{array}\!\!\right) 
					  \sim \left(\!\!\begin{array}{cc|c}4 & 4 & 7 \\ 0 & 0 & 0\end{array}\!\!\right) 
					  \sim \left(\!\!\begin{array}{cc|c}1 & 1 & 7/4 \\ 0 & 0 & 0\end{array}\!\!\right).\]
					We find that \[\begin{bmatrix}\alpha \\ \beta\end{bmatrix} = \begin{bmatrix}7/4 \\ 0\end{bmatrix} 
					+ \beta\begin{bmatrix}-1 \\ 1\end{bmatrix},\] where \(\beta\) is the free variable. Each choice of \(\alpha, \beta\) obtained from the
					parametric vector equation yields a curve with the same amount of error.

				\item \[R\vec{\hat{x}} = Q^T\vec{b} \implies \begin{bmatrix}2 & 2 \\ 0 & 1\end{bmatrix}\vec{\hat{x}}
								 = \sqrt{2} \cdot \frac{1}{\sqrt{2}} \begin{bmatrix}1 & - 1 \\ 1 & 1\end{bmatrix}
								 \begin{bmatrix}2 \\ 1\end{bmatrix} = \begin{bmatrix}1 \\ 3\end{bmatrix}.\]

					Row-reduce the augmented matrix to reach to the solution. 
					\[\left(\!\!\begin{array}{cc|c}2 & 2 & 1 \\ 0 & 1 & 3\end{array}\!\!\right) 
					  \sim \left(\!\!\begin{array}{cc|c}2 & 0 & -5 \\ 0 & 1 & 3\end{array}\!\!\right)
					  \sim \left(\!\!\begin{array}{cc|c}1 & 0 & -5/2 \\ 0 & 1 & 3\end{array}\!\!\right).\]
					We get that \(\vec{\hat{x}} = \begin{bmatrix}-5/2 \\ 3\end{bmatrix}\).

				\item Since the columns were already orthogonal, the solution should just be the weights used in the projection of \(\vec{b}\) onto the
					the column space of \(A\). 
					\[\hat{x}_1 = \frac{\vec{a}_1 \cdot \vec{b}}{\vec{a}_1 \cdot \vec{a}_1} = -\frac{3}{2}, \quad 
					   \hat{x}_2 = \frac{\vec{a}_2 \cdot \vec{b}}{\vec{a}_2 \cdot \vec{a}_2} =\frac{5}{2}\]
			\end{enumerate}
	\end{enumerate}

\end{document}
